\documentclass[aspectratio=169]{ctexbeamer}
\usepackage[T1]{fontenc}
\usepackage{latexsym,amsmath,xcolor,multicol,booktabs,calligra}
\usepackage{graphicx,listings,stackengine}
\usepackage{hyperref}

\usepackage{PekingU}

% 去掉每节/子节前面的自动目录页
\AtBeginSection[]{}
\AtBeginSubsection[]{}

% ===== 每节课改这里 =====
\author{Cap1taL}
\title{贪心与搜索}
\date{2026年7月}
% ========================

% 代码高亮设置
\definecolor{deepblue}{rgb}{0,0,0.5}
\definecolor{deepred}{rgb}{0.6,0,0}
\definecolor{deepgreen}{rgb}{0,0.5,0}
\definecolor{halfgray}{gray}{0.55}

\lstset{
    basicstyle=\ttfamily\small,
    keywordstyle=\bfseries\color{deepblue},
    emphstyle=\ttfamily\color{deepred},
    stringstyle=\color{deepgreen},
    numbers=left,
    numberstyle=\small\color{halfgray},
    rulesepcolor=\color{red!20!green!20!blue!20},
    frame=shadowbox,
}

\begin{document}

\kaishu

% 封面
\begin{frame}
    \titlepage
\end{frame}

% 目录
% \begin{frame}{目录}
%     \tableofcontents[sectionstyle=show,subsectionstyle=show/shaded/hide,subsubsectionstyle=show/shaded/hide]
% \end{frame}

% ===== 从这里开始写你的内容 =====

\section{题目}


% A
\subsection{CodeForces 496E Distributing Parts}
\begin{frame}{CodeForces 496E Distributing Parts}
    热身题 \\

    $n$ 首歌曲，每首的音域为 $[a_i,b_i]$

    $m$ 个演员，每个人的音域范围为 $[c_i,d_i]$，他最多可以演唱 $k_i$ 首歌曲

    一首歌曲能被一个演员演唱，当且仅当演员 $[c,d]$ 完全包含歌曲 $[a,b]$
    
    判断是否存在一个分配歌曲演唱的方案，并给出构造

    $n,m\leq 10^5$

    其余不重要
\end{frame}

\begin{frame}{Solution}
    先忽略 $k$ 的影响，不难转化为经典模型
    
    将演员与歌曲，一同按区间左端点升序排列，从左向右依次考虑它们
    
    \begin{enumerate}
        \item 若考虑到演员-左端，将其纳入维护
        \item 若考虑到演员-右端，将其移出维护
        \item 若考虑到歌曲-左端，我们从维护的演员里，找到\textcolor{peaking}{能演唱}当前歌曲，且\textcolor{peaking}{右端最小}的的演员，将其与当前乐曲匹配
    \end{enumerate}
\end{frame}

\begin{frame}{Solution}
    现在考虑 $k$ 的影响
    
    我们将一个人可唱 $k$ 次，等效为 $k$ 个人，则可直接套用刚才的做法
    
    $k$ 可以到 $10^9$ 量级，因此我们不能真的拆人（？！）
    
    我们在纳入维护演员时，为每个人额外维护一个“剩余次数”即可
    
    具体维护可以使用 set 或者 multiset
\end{frame}

% B
\begin{frame}{CodeForces 525D Arthur and Walls}
    题目大意
\end{frame}

% C
\begin{frame}{CodeForces 442C Artem and Array}
    题目大意
\end{frame}

% D
\begin{frame}{洛谷 P6243 The Milk Queue G}
    题目大意
\end{frame}

% E
\begin{frame}{洛谷 P6631 序列}
    题目大意
\end{frame}

% F
\begin{frame}{洛谷 P8314 Parkovi}
    题目大意
\end{frame}

% G
\begin{frame}{洛谷 P5665 划分}
    题目大意
\end{frame}

% H
\begin{frame}{QOJ 5173 染色}
    题目大意
\end{frame}

% I
\begin{frame}{QOJ 1173 Knowledge Is...}
    题目大意
\end{frame}

% J
\begin{frame}{CodeForces 802M3 April Fools' Problem (hard)}
    题目大意
\end{frame}

% K
\begin{frame}{CodeForces 436E Cardboard Box}
    题目大意
\end{frame}

% L
\begin{frame}{洛谷 P4511 日程管理}
    题目大意
\end{frame}

% M
\begin{frame}{SPOJ MTREECOL Color a tree}
    题目大意
\end{frame}

% N
\begin{frame}{洛谷 P3620 数据备份}
    题目大意
\end{frame}

% O
\begin{frame}{CodeForces 2128F Strict Triangle}
    题目大意
\end{frame}

% P
\begin{frame}{CodeForces 2034G2 Simurgh's Watch (Hard Version)}
    题目大意
\end{frame}

% ===== 结束页 =====

\begin{frame}
    \begin{center}
        {\Huge\calligra Thanks!}
    \end{center}
\end{frame}

\end{document}
.
